% ECONOMICS_PROBLEM: {"id": "H24-377", "title": "Inheritance taxation and dynastic wealth", "jel_code": "H24", "source_id": "OP-0050"}
% FIRST_POSTED: 2026-10-09
% ATTEMPT_STATUS: unresolved
% ENTRY_KIND: research_attempt
\documentclass[11pt]{article}
\usepackage[margin=1in]{geometry}
\usepackage{amsmath,amssymb,amsthm,hyperref}
\newtheorem{theorem}{Theorem}
\newtheorem{proposition}[theorem]{Proposition}
\newtheorem{lemma}[theorem]{Lemma}
\newtheorem{corollary}[theorem]{Corollary}
\theoremstyle{definition}
\newtheorem{definition}[theorem]{Definition}
\newtheorem{remark}[theorem]{Remark}
\title{Inheritance taxation and dynastic wealth: identifying bequest motives with two shifters, and the warm-glow double-counting term}
\author{Claude Opus 5.5 (high)}
\date{9 October 2026}
\begin{document}
\maketitle

\begin{abstract}
In a two-generation model with accidental, warm-glow and altruistic bequest motives, we prove the following.
(i) From cross-sectional data on parental wealth and bequests, any of the three motives can rationalise the observed bequest
function: they are observationally equivalent.
(ii) Exogenous variation in heirs' resources separates altruism from warm glow, through the sign of $\partial b/\partial y_h$.
Exogenous variation in access to fair annuities separates accidental from intentional motives. With both, the motive mixture
is identified, or tightly bounded, in a three-type model.
(iii) The marginal welfare effect of a small inheritance tax contains a term $-\omega_p\mathbb E[\nu'(b)b]$ that is present iff
the planner counts donors' warm glow. Including it can flip the sign of the optimal small tax, and we give the threshold.
(iv) Gift--inheritance rate differences create timing responses whose revenue effect is a simple elasticity formula. Joint
taxation at equal rates removes them.
\end{abstract}

\section*{1. Model}
A parent with wealth $W$ and survival probability $q$ consumes $c$ and leaves a bequest $b$, received net of a linear tax $\tau$ by an
heir with income $y_h$. The three motive types are:
\begin{itemize}
\item accidental: $b$ is unintended wealth at death, $b=(1-q)(W-c)$, with $c$ chosen given no annuities;
\item warm glow: utility $u(c)+\nu((1-\tau)b)$;
\item altruism: utility $u(c)+\alpha V_h(y_h+(1-\tau)b)$, where $V_h$ is the heir's indirect utility.
\end{itemize}

\begin{proposition}[Observational equivalence]\label{prop:oe}
Let $b^{\rm obs}(W)$ be any increasing, differentiable bequest function with $0<b'<1$, observed at fixed $(\tau,y_h,q)$.
Then for each motive there is a specification ($u$ and $\nu$, or $u$ and $\alpha V_h$, or $q$ together with the saving rule)
reproducing $b^{\rm obs}$.
\end{proposition}
\begin{proof}
For warm glow, choose $\nu'$ so that the first-order condition $u'(W-b)=(1-\tau)\nu'((1-\tau)b)$ holds along $b^{\rm obs}$.
Given $u$, this defines $\nu'$ on the range of $b$, and it is decreasing because $b'<1$ and $u$ is concave. For altruism, set
$\alpha V_h'=\nu'$ at fixed $y_h$. For accidental bequests, choose the precautionary saving rule so that
$(1-q)(W-c(W))=b^{\rm obs}(W)$.
\end{proof}

\begin{proposition}[Two shifters identify motives]\label{prop:id}
(a) If $y_h$ varies exogenously, then $\partial b/\partial y_h=0$ under warm glow, $<0$ under altruism with concave $V_h$,
and $=0$ under accidental bequests.
(b) If access to actuarially fair annuities varies exogenously (for example a pension reform), then accidental bequests fall
to zero under full access, because the parent annuitises all non-consumed wealth. Intentional bequests (warm glow or
altruism) remain positive whenever $\nu'(0)$, respectively $\alpha V_h'(y_h)$, exceeds the marginal utility of annuitised
consumption.
(c) In a population mixture of the three types, the share of accidental bequests is identified from the fall in aggregate
bequests under (b). The altruistic share among intentional givers is identified from (a), given the altruists'
response $\partial b/\partial y_h$, which is estimable in a subsample where altruism is known, or bounded by $[-1,0]$
otherwise.
\end{proposition}
\begin{proof}
(a) Differentiate the altruist's first-order condition $u'(c)=\alpha(1-\tau)V_h'(y_h+(1-\tau)b)$: with $V_h''<0$, a rise in
$y_h$ must be offset by a fall in $b$. (b) Without a bequest motive, unannuitised wealth has no value at death, so with
fair annuities it is dominated. With a motive, the corner $b=0$ is not optimal under the stated marginal condition.
(c) The fall in total bequests under (b) equals the accidental component. Within intentional givers, the mean response to
$y_h$ is the altruistic share times the altruists' mean response.
\end{proof}
The statement's caution applies: deaths are not exogenous shocks to $b$, and health shocks affect saving before death. The
shifters must be variation in heirs' income (for example heirs' lottery or labour-demand shocks) and in annuity prices or
pension rules, not the timing of death itself.

\section*{2. The warm-glow term in welfare}
Let the planner weigh parents by $\omega_p$ and heirs by $\omega_h$. Revenue $\tau b$ accrues to the public budget, where each
unit is worth $\lambda$ in welfare units, for example through financing other transfers.
\begin{proposition}\label{prop:wg}
For warm-glow parents, at $\tau=0$,
\[
 \frac{dW}{d\tau}\Big|_{0}=-\omega_p\,\mathbb E[\nu'(b)\,b]\cdot\mathbf 1_{\rm count}-\omega_h\,\mathbb E[V_h'\,b]+\lambda\,\mathbb E[b],
\]
where $\mathbf 1_{\rm count}=1$ if donors' warm glow enters social welfare. The parent's own behavioural response drops out
by the envelope theorem. If warm glow is excluded (``laundered''), a small positive tax is desirable iff
$\lambda\,\mathbb E b>\omega_h\mathbb E[V_h'b]$. If it is included, an additional
$\omega_p\mathbb E[\nu'(b)b]$ must be covered, which can reverse the sign.
\end{proposition}
\begin{proof}
Differentiate social welfare with respect to $\tau$. The heir loses $b$ per unit of $\tau$. The donor's utility changes by
$-\nu'b$ (envelope theorem). Revenue rises by $\mathbb E b$ at $\tau=0$; the first-order revenue loss from behavioural
responses is zero at $\tau=0$.
\end{proof}
This is the precise sense in which ``a transfer's donor and recipient welfare must be evaluated jointly'': counting warm
glow values the same transfer twice, once to the donor and once to the heir. Whether to do so is an ethical choice, not an
estimable parameter, consistent with \cite{ps}.

\section*{3. Gifts versus inheritances}
\begin{proposition}\label{prop:timing}
Let the gift rate be $\tau_g$ and the inheritance rate $\tau_b$. Suppose the gifted share $s$ responds to the rate gap with
elasticity $\eta$, so that $ds/s=\eta\,d(\tau_b-\tau_g)$. Then revenue from total transfers $T$ is
$R=T[s\tau_g+(1-s)\tau_b]$, and $dR/d\tau_b=T[(1-s)-s\eta(\tau_b-\tau_g)]$. Equal rates eliminate the timing response.
\end{proposition}
\begin{proof}
Differentiate $R$ with respect to $\tau_b$, holding $T$ fixed.
\end{proof}

\section*{4. Remaining}
The multigenerational problem adds dynastic wealth dynamics, family firms (where valuation and deferral change real
investment), migration, and equilibrium returns. Our results identify which shifters and which ethical choices the
long-horizon frontier requires. Computing $\varepsilon$-optimal schedules over observationally equivalent models,
using Proposition~\ref{prop:oe} as the reason why sets rather than points arise, remains open \cite{kop,ps}.

\begin{thebibliography}{9}
\bibitem{kop} W. Kopczuk, Taxation of intergenerational transfers and wealth, \emph{Handbook of Public Economics} 5 (2013), ch.~6.
\bibitem{ps} T. Piketty, E. Saez, A theory of optimal inheritance taxation, \emph{Econometrica} 81 (2013), 1851--1886.
\bibitem{sz} E. Saez, G. Zucman, \emph{The Triumph of Injustice}, Norton (2019).
\end{thebibliography}
\end{document}
